\section{Residues at primitive-recursively named singularities}\label{sec:named-residue}

The repository's \src{docs/proofs/analysis/cauchy_residue_ipra/report.tex}{Cauchy and residue manuscript} proves a quantitative
residue theorem for rational polygonal fillings and explicitly restricts the
singularity centres to rational points. A rational point close to a singularity
generally has residue zero. Using it as the centre of a Laurent expansion would
therefore give a wrong answer. A rational point can instead locate the
\emph{hole in a triangle chain}, while a translated chart computes the Laurent
expansion about the actual singularity. We prove this change of presentation.
The base analytic results we reuse are the finite-chain Cauchy theorem, square
kernel moment and annular Laurent estimate of that manuscript. These are inputs
from its human proofs. Their deep Lean formalizations are dependencies of the
proposed formalization here.

\subsection{Quantitative chart transport}

Write $|z|_1=|\Re z|+|\Im z|$, $|z|_\infty=\max(|\Re z|,|\Im z|)$ and
$Q_r(b)=\{z:|z-b|_\infty\le r\}$. Our centres are given by a fixed PR
procedure $A(t)\in\mathbb Q(i)$ satisfying
$|A(t)-A(u)|_1\le2^{-t}+2^{-u}$. Write $a$ for its name. For a finite
family of variable size, use a single procedure $A(j,t)$, and likewise
single indexed PR procedures for the evaluators, bounds and moduli on
the local carriers.
Whenever a named quantity occurs in an estimate it denotes its regular
rational approximation inequalities, with the errors of its samples
added. For example, $|a-A(t)|_1\le2^{-t}$ means the displayed
two-index inequality.

On a rational carrier give uniform PR evaluators $F(x,n),D(x,n)$ for a
function $f$ and its derivative coefficient $d$, each of accuracy $2^{-n}$,
and magnitude bounds. Precisely, the samples of each at a fixed rational
$x$ satisfy the same two-index inequality as $A$. For rational $x,y$ in
the carrier, require the following uniform bounds with supplied PR moduli $\nu,\mu$.
\begin{align*}
 |y-x|_1\le2^{-\nu(k)}&\ \Longrightarrow\ |f(y)-f(x)|_1\le2^{-k},\\
 |y-x|_1\le2^{-\mu(k)}&\ \Longrightarrow\
 |f(y)-f(x)-d(x)(y-x)|_1\le2^{-k}|y-x|_1 .
\end{align*}
The finite reading of the second bound at common sample precision $n$
adds $(2+|y-x|_1)2^{-n}$ to the right. These conditions compare the
indicated pairs of points, and do not require a path between them.
All moduli may be made nondecreasing by taking a running maximum.

Fix a rational $r>0$, choose $s_0$ with $2^{-s_0}\le r/128$, and set
$b=A(s_0)$. For instance $s_0=\lceil128/r\rceil$ suffices. Assume the
above data on the closed rational annulus
\[
 K=\{x:r/32\le |x-b|_\infty\le9r\}.
\]
This is an explicit margin assumption. A presentation supplied on all
separated compact subannuli about $a$ need only supply its data here, since
$K$ lies inside $r/64\le|x-a|_\infty\le10r$. For $t\ge s_0$,
$|A(t)-b|_1\le r/64$. Thus $A(t)+w$, for
$r/8\le |w|_\infty\le8r$, remains in $K$ even after an additional
taxicab displacement of $r/64$.

\begin{lemma}[Translated evaluator on the closed annulus]\label{lem:analytic-shift}
The function $g(w)=f(a+w)$, on rational $w$ with
$r/8\le |w|_\infty\le8r$, has PR evaluators, magnitude bounds,
and the same continuity and first-order remainder moduli. A modulus for
variation of the original derivative on the sampled points is derivable
from $\nu,\mu$ and the margin.
\end{lemma}
\begin{proof}
Given $k$, let $H=\max\{\mu(k+2),\lceil64/r\rceil\}$ and
$h=2^{-H}$, a positive \emph{real} rational step. At every sampled point
$x=A(t)+w$, the points $x$ and $x+h$ belong to $K$. The secant
$[f(x+h)-f(x)]/h$ differs from $d(x)$ by at most $2^{-k-2}$. If two
such points $x,y$ satisfy $|x-y|_1\le2^{-\nu(k+H+2)}$, continuity at
$x,y$ and at $x+h,y+h$ bounds the difference of the secants by
$2\cdot2^{-k-H-2}/h=2^{-k-1}$. The two end errors give the derivative
modulus $\nu_D(k)=\nu(k+H+2)$. The division is only by the explicitly
chosen rational $h$. Comparing secants made of finite samples adds at
most $4\cdot 2^{-n}/h$; a sample precision $n\ge q+H+2$ bounds this
extra term by $2^{-q}$. This also justifies the regular-name shorthand
used in the calculation.

Define
\[
 e(n)=\max\{s_0,n+3,\nu(n+3)+2,\nu_D(n+3)+2\}.
\]
Evaluate the new function and derivative at $(w,n)$ by applying $F,D$
at $(A(e(n))+w,n+3)$. For any $t\ge e(n)$ the two rational arguments
$A(t)+w$ and $A(e(n))+w$ differ by at most $2^{1-e(n)}$. The
continuity bounds just obtained, plus the error of the returned sample,
put it within $2^{-n-2}$ of the corresponding value at $A(t)+w$.
Taking a common $t$ for two proposed samples proves regularity, for
both procedures. To compare $w$ and $v$, choose the same $t$ in both
translates. Their difference is exactly $v-w$, so both input bounds
apply with their original moduli. Moving the function and derivative
samples back costs at most their prescribed sample errors. In particular
the new remainder slack at common precision $n$ is bounded by
$(2+|v-w|_1)2^{-n}$. The argument uses only arbitrarily accurate
rational samples; clearing any auxiliary $2^{-q}$ error in a rational
inequality uses decidable comparison and the denominator of a putative
positive gap. Magnitude bounds transfer in the same way.
\end{proof}

\paragraph{The finite analytic input in more detail.}
Here is a reconstruction of the estimates behind the three inherited facts. For a rational edge $[p,q]$ set $z(t)=p+t(q-p)$ and use left Riemann sums on a dyadic subdivision. If the integrand varies by at most $2^{-k}$ on each panel, comparing a sum with any dyadic refinement costs at most $|q-p|_1 2^{-k}$, plus the length times the sum of the two point-evaluation errors for independently sampled sums. For a continuity modulus $\nu$, taking $2^m\ge |q-p|_1 2^{\nu(k)}$ ensures the required panel size. Sorting the union of two rational node lists gives a common refinement, also for a split at a nondyadic rational parameter, and the same oscillation bound applies. This proves additivity and orientation reversal as name equalities.

For a triangle with taxicab diameter $D$ and perimeter $P$, subdivide repeatedly by the three midpoints. At level $m$ there are $4^m$ triangles of diameter at most $D2^{-m}$ and perimeter $P2^{-m}$. Choose this diameter below the input Taylor radius at $k$. On each triangle subtract the affine function $f(x)+d(x)(z-x)$ at a fixed vertex $x$. Its integral on the triangle boundary is zero, by integrating the constant and linear terms along the three edges, or by telescoping their quadratic primitive. The remaining integrand has norm at most $D2^{-m}2^{-k}$. All internal edges cancel, and the total integral has norm at most $DP2^{-k}$. For finite sampled sums the total perimeter at this level is $2^mP$, so extra error $\varepsilon$ is achieved by per-unit-length error $\varepsilon/(1+2^mP)$. This proves the finite-chain Cauchy identity by an arbitrary rational tolerance, including signed multiplicities.

For the square kernel, define $\pi=4\int_0^1(1+t^2)^{-1}\,dt$; the integral has a PR modulus from a rational Lipschitz bound and $\pi\ge2$. On the right edge of the unit square the integrand $dz/z$ is $(t+i)\,dt/(1+t^2)$ for $-1\le t\le1$. The odd part cancels and rotation gives the same value on all four sides. Hence the total is $2\pi i$, at any positive scale by change of variables. For $z^j$ with integer $j\ne-1$, $z^{j+1}/(j+1)$ is a primitive on the contour; the usual telescoping edge proof follows from the finite binomial expansion for positive powers, and from the difference of inverses under the edge's positive separation for negative ones.

For completeness, the annular Cauchy formula at a rational $v\in\partial Q_r(0)$ uses the kernel $g(u)/(u-v)$. Exclude a rational square of half-side $h$ about $v$ lying between $Q_s$ and $Q_R$. The remaining square region with two disjoint square holes is triangulated by drawing horizontal and vertical lines through all sides, retaining its rectangles and taking two triangles in each. Division by $u-v$ there is separated: multiplication and the inverse difference identity supply the required moduli. The Cauchy identity equates outer minus inner integrals to the integral on the small square. Replacing $g(u)$ by $g(v)$ there has integral error at most $(8h)(2/h)2^{-k}$ once $2h\le2^{-\nu(k)}$. The square moment then gives $2\pi i g(v)$ to arbitrary precision. Finally expand the two kernels by the \emph{finite} geometric identity with remainder (on the inner square expand in $u/v$, on the outer square in $v/u$). The explicit bound below results from their $1/2$ ratios. Thus the inherited analytic argument uses rational quadrature and the stated modulus of differentiability, even when $v$ is on an intermediate square edge. These proofs still require deep formalization for their use in ordinary iPRA.


\subsection{From local Laurent coefficients to a global integral}

\begin{theorem}[Named-centre residue formula]\label{thm:analytic-named}
Let $\Gamma$ be a finite closed rational polygonal integer chain and
let $a_j$ be finitely many named points with rational data $r_j,s_j,b_j,K_j$
as above. The rational holes $Q_{r_j}(b_j)$ are pairwise disjoint and
avoid $\Gamma$. Give effective holomorphy data on each $K_j$, and on
a rational triangle chain $T$ lying outside all open holes, such that
\[
 \partial T=\Gamma-\sum_j w_j\,\partial Q_{r_j}(b_j)
\]
with equality after common edge subdivision. The function evaluations
agree on overlaps. Then the Laurent coefficient $c_{j,-1}$ of $f$ in
the coordinate $z-a_j$ has a PR name and
\[
 \int_\Gamma f(z)\,dz=2\pi i\sum_j w_j c_{j,-1},
 \qquad w_j=(2\pi i)^{-1}\int_\Gamma\frac{dz}{z-a_j}.
\]
The coefficient is local to the stated annulus. When the data extend
coherently to a punctured neighbourhood, it is the residue of that germ.
There is no bound on pole order in the hypotheses.
\end{theorem}
\begin{proof}
For one hole, use Lemma~\ref{lem:analytic-shift} to construct $g$.
Apply the rational-centre annular theorem at the rational point $0$
in the $w$-coordinate. To specify the computation, put $R=8r$ and
$s=r/8$. For $n\ge0$ its coefficients are
\begin{align*}
 c_n &=(2\pi i)^{-1}\int_{\partial Q_R(0)}g(w)w^{-n-1}\,dw,\\
 c_{-n-1}&=(2\pi i)^{-1}\int_{\partial Q_s(0)}g(w)w^n\,dw .
\end{align*}
With outer and inner magnitude bounds $M_R,M_s$, its proof bounds the
error of the sum from $-N$ through $N-1$ on $\partial Q_r(0)$ by
$C2^{-N}$, where
\[
 C=16RM_R/(R-r)+16sM_s/(r-s).
\]
Indeed the ratios $|w/u|_1$ at the outer boundary and $|u/w|_1$ at the
inner boundary are each at most $1/2$ in the finite geometric expansions
of the Cauchy kernels. The square moments vanish in all degrees other than
$-1$ and equal $2\pi i$ in that degree. Integrating the approximation
therefore yields $\int_{\partial Q_r(0)}g=2\pi i c_{-1}$ with error
bounded by $8rC2^{-N}$ before passage to names. The same moment test
after multiplication by $w^{-m-1}$ recovers each coefficient of any
PR Laurent presentation with a supplied uniform tail modulus on the square, establishing
its uniqueness. Thus the coefficient refers to the specified centre.

It remains to identify this local integral with the boundary integral
used in the rational chain. For each $t\ge s_0$ join corresponding
vertices of $\partial Q_r(b)$ and $\partial Q_r(A(t))$ and triangulate
each of the four panels by one diagonal. A point in a panel is a point
on the corresponding side of $Q_r(b)$ displaced by a scalar multiple
between zero and one of $A(t)-b$. Hence its max-norm distance from $b$
is between $r-r/64$ and $r+r/64$. All eight rational triangles lie in
$K$, including degenerate panels. Orient them in the order
$[p,q,q'],[p,q',p']$ for bottom side $p q$ and top side $p'q'$.
Diagonals and connectors cancel, leaving the bottom boundary minus
the top boundary. The finite-chain Cauchy theorem consequently gives
identical integrals on the two squares. Rational translation on each
edge changes the latter integral to
$\int_{\partial Q_r(0)} f(A(t)+w)\,dw$ by literal equality of
parametrized quadrature sums. For $t\ge\max(s_0,\nu(k)+2)$ the
integrand differs uniformly from $g$ by at most $2^{-k}$, by comparison
with all later rational translates. The integral error is at most
$8r2^{-k}$ by the perimeter bound on quadrature. As $k$ is arbitrary
the proxy integral equals $2\pi i c_{-1}$. Applying the finite-chain
Cauchy theorem on $T$ and summing these identities proves the first
claim. For example, to make the total translation error $\le2^{-q}$,
take a common $k=q+\lceil1+\sum_j8r_j|w_j|\rceil$. This gives a
PR schedule also when the number of holes varies.

For the second identity fix $v=A_j(t)$ with $t\ge s_j$. On $T$ the
kernel $(z-v)^{-1}$ has positive separation because $T$ avoids the open
$j$th hole and $|v-b_j|_1\le r_j/64$. Its integral around that hole
equals $2\pi i$ by the same rational eight-triangle construction and
the centred kernel moment. On a different hole the kernel extends over
its entire square; the two triangles filling that square give zero.
The chain identity thus gives the normalized integral $w_j$.
Put $L$ equal to the total taxicab length of $\Gamma$, counting
absolute multiplicities, and $\delta=r_j/2$. Both centre distances
along $\Gamma$ are at least $\delta$ in max norm. Since
$|u^{-1}|_1\le2/\delta$ when $|u|_\infty\ge\delta$, the inverse
difference identity bounds the integral change on replacing $v$ by
$a_j$ by $4L2^{-t}/\delta^2$. Taking arbitrary precision proves the
second identity and shows independence of the proxy for the index.
Radius independence for a punctured germ follows by applying the
centred Cauchy theorem to the intervening annulus in the translated
coordinate. All constructions use bounded lists of rational
operations, the given moduli and precision shifts.
\end{proof}

The geometric hypotheses are attainable without knowing an exact
centre. For example, given a rational lower bound $d>0$ on all
pairwise centre and centre-to-contour max-norm separations, choose a
common $r\le d/32$ small enough for the analytic collars and approximate
each centre to $r/128$. The rational holes then satisfy the required
separations by the triangle inequality. The repository's rational
planar filling algorithm applies to their proxies. For a restricted
analytic domain the membership of its retained triangles still requires
evidence. In the proof the singularities may have nonrational real and
imaginary parts. The information consumed about them consists of
approximations at finitely prescribed precisions and the supplied
separations.

\subsection{An irrational double pole and its exact integral}

Take $a=\sqrt2+i\sqrt3$, $f(z)=z^2/(z-a)^2$, and let $\Gamma$ run
counterclockwise around $[1,2]^2$. A rational approximation algorithm
uses $m=2^{t+2}$ and returns
\[
 A(t)=m^{-1}\bigl(\lfloor\sqrt{2m^2}\rfloor
                 +i\lfloor\sqrt{3m^2}\rfloor\bigr).
\]
The integer square roots are bounded searches for the largest $u\le N$
with $u^2\le N$. The resulting two coordinates have error less than
$1/m$ each. Choose $r=1/64$, $s_0=14$, so $2^{-s_0}\le r/128$.
The rational comparisons $7/5<\sqrt2<3/2$ and
$17/10<\sqrt3<7/4$ follow on squaring. They put the whole $9r$
proxy square inside $[1,2]^2$. The horizontal and vertical lines
through the proxy hole sides divide the outer square into nine
rectangles; the eight outside the hole give sixteen triangles and
$w=1$.

Here are witnesses for the analytic hypotheses. The distance to $a$ in
max norm on $T\cup K$ is at least $r/64$. With $U=256/r$, the inverse
$u=(z-a)^{-1}$ and all its sufficiently precise rational approximations
are bounded in taxicab norm by $U$: evaluate at $A(t)$ with
$2^{-t}\le r/128$ and use the inverse bound. The inverse difference
identity gives $|u(y)-u(x)|_1\le U^2|y-x|_1$, and the first-order
remainder for $u$ has bound $U^3|y-x|_1^2$, since
\[
 u(y)-u(x)+u(x)^2(y-x)=u(x)^2u(y)(y-x)^2 .
\]
Expansion of $u(y)^2$ at $x$ bounds the corresponding remainder for
$u^2$ by $3U^4|y-x|_1^2$. Indeed the two terms are twice $u(x)$
times the preceding remainder and $(u(y)-u(x))^2$. Since
$f=1+2au+a^2u^2$ and $|a|_1\le4$, continuity is Lipschitz with constant
$8U^2+32U^3$, and its quadratic remainder constant is $8U^3+48U^4$.
For any such constant $B$, $k+\lceil B+1\rceil$ is a sufficient dyadic
modulus. Magnitude bounds follow from the same expression for $f$
and $d=-2au^2-2a^2u^3$. To evaluate, use the equivalent formulas
$f=z^2u^2$ and $d=2zu^2-2z^2u^3$. Since $|z|_1\le4$ on the
carriers, replacing $a$ by $A(t)$ in $u$ incurs errors at most
$32U^3 2^{-t}$ for $f$ and $(16U^3+96U^4)2^{-t}$ for $d$, by factoring
the differences of the second and third powers. This supplies fixed
PR evaluators with the required service precisions.

Finally $f(a+w)=1+2a/w+a^2/w^2$ by multiplication by $w^2$.
The exact integral is consequently
\[
 \int_\Gamma \frac{z^2}{(z-a)^2}\,dz=4\pi i(\sqrt2+i\sqrt3).
\]
In particular, the imaginary part of the centre affects the real
part of the integral. The rational proxy itself is a regular point of $f$.

\subsection{Lean targets for this extension}

The included \texttt{Secant.lean} proves the two-component rational
estimate underlying Lemma~\ref{lem:analytic-shift}. If the two
first-difference residuals are $\le R$ and the two cross-differences
are $\le u$, it proves $h|d-e|_1\le2R+2u$. Substituting
$R=h2^{-k-2}$ and $u=2^{-k-H-2}$ gives the stated modulus.
Its proof uses Mathlib's rational arithmetic at Lean 4.32.0.
The included \texttt{ContourRibbon.lean} proves
\texttt{panel\_cancel}, \texttt{strip\_cancel}, and \texttt{closed\_strip}
in Lean's ordinary logic. For any antisymmetric integer observation of
a directed edge, the signed sum over the two triangles of each panel is
the bottom edge minus the top edge plus its end connector minus its
start connector. Induction over consecutive pairs then leaves only
the final and initial connectors, which cancel for a loop. Taking the
observation to be the coefficient of one queried oriented edge gives
the required identity of chains. The verifier uses Lean 4.32.0 and
\texttt{Std}. To internalize the remaining argument:
\begin{enumerate}
\item Use integer numerator and positive denominator codes; implement
      componentwise translation, the two norms, inverse at separated
      points, and the predicates for $K$ and the core annulus by
      rational cross multiplication. Define the strip as the two listed
      triangles per paired edge. Its object induction invariant, with
      the queried edge fixed as parameter, is the integer equation of
      \texttt{strip\_cancel}.
\item For fixed codes for $A,F,D,\nu,\mu$, express the two input
      inequalities on rational samples, plus their coherence, as local
      hypotheses of \texttt{Prf PRAAx}. Define $H,e$ by composition.
      Prove the secant estimate at a common auxiliary precision and use
      the positive denominator of a rational gap to discharge extra
      slack. Apply the input inequalities at a common translate to
      derive the two new contracts. Universally close over parameters
      after deriving these finite inequalities.
\item Import or implement the manuscript's finite quadrature and
      Cauchy estimates, square moments and Laurent cutoff. Instantiate
      them with the translated evaluators and rational centre zero.
      For the eight triangle applications use the rational band
      inclusion proved above. For a variable family fold weighted
      integral discrepancies, with invariant that their norm is at
      most the partial sum of the individual rational error budgets.
      The centre-approximation argument for the kernel uses only
      multiplication of separated inverses and a scalar perimeter
      bound. These steps target the finite precision readings of both
      identities in Theorem~\ref{thm:analytic-named}.
\end{enumerate}
